\Needspace{12\baselineskip}
\begin{center}
\textit{Formen mit der Determinante 0.}
\end{center}

\Needspace{8\baselineskip}
\Article{215}
Bisher haben wir Formen mit der Determinante 0 aus allen Untersuchungen ausgeschlossen. Damit unsere Theorie in jeder Hinsicht vollständig wird, müssen wir noch einiges über sie hinzufügen. Wir haben allgemein bewiesen: Enthält eine Form mit der Determinante $D$ eine Form mit der Determinante $D'$, so ist $D'$ ein Vielfaches von $D$. Daher kann eine Form mit der Determinante $=0$ offenbar nur eine Form enthalten, deren Determinante ebenfalls $=0$ ist. Es bleiben also nur zwei Aufgaben: $1^{\mathrm{o}}$ \textit{Für zwei gegebene Formen $f,F$, von denen die zweite die Determinante 0 hat, entscheiden, ob die erste die zweite enthält, und in diesem Fall alle Transformationen der ersten in die zweite angeben.} $2^{\mathrm{o}}$ \textit{Alle Darstellungen einer gegebenen Zahl durch eine gegebene Form mit der Determinante 0 finden.} Die erste Aufgabe erfordert eine andere
\SourceLeaf{217}
Methode, wenn die Determinante der ersten Form $f$ ebenfalls 0 ist, als wenn sie nicht 0 ist. Dies alles behandeln wir jetzt.

I. Zunächst bemerken wir, dass sich jede Form $axx+2bxy+cyy$ mit der Determinante $bb-ac=0$ als $m(gx+hy)^2$ schreiben lässt, wobei $g,h$ teilerfremd sind und $m$ eine ganze Zahl ist. Sei nämlich $m$ der größte gemeinsame Teiler von $a,c$, mit demselben Vorzeichen genommen wie diese Zahlen selbst (dass sie keine entgegengesetzten Vorzeichen haben können, sehen wir leicht). Dann sind $\frac{a}{m},\frac{c}{m}$ teilerfremde nichtnegative ganze Zahlen. Ihr Produkt ist $=\frac{bb}{mm}$, \textit{d.h.} eine Quadratzahl, und daher sind sie selbst Quadratzahlen (Artikel~21). Sei $\frac{a}{m}=gg$, $\frac{c}{m}=hh$. Dann sind auch $g,h$ teilerfremd; außerdem gilt $gghh=\frac{bb}{mm}$ und $gh=\pm\frac{b}{m}$. Daraus folgt offenbar
\[
m(gx\pm hy)^2=axx+2bxy+cyy
\]

Nun seien zwei Formen $f,F$ gegeben, beide mit der Determinante 0, und es sei
\[
f=m(gx+hy)^2,\quad F=M(GX+HY)^2
\]
wobei $g$ zu $h$ und $G$ zu $H$ teilerfremd ist. Ich behaupte: Enthält $f$ die Form $F$, so ist $m$ entweder gleich $M$ oder teilt zumindest $M$, und der Quotient ist eine Quadratzahl. Ist umgekehrt $\frac{M}{m}$ eine ganze Quadratzahl, so ist $F$ unter $f$ enthalten. Nehmen wir nämlich an, $f$ gehe durch die Substitution
\[
x=\alpha X+\beta Y,\quad y=\gamma X+\delta Y
\]
in $F$ über. Dann ist
\[
\frac{M}{m}(GX+HY)^2=((\alpha g+\gamma h)X+(\beta g+\delta h)Y)^2
\]
woraus leicht folgt, dass $\frac{M}{m}$ eine Quadratzahl ist. Wir setzen sie $=ee$ und erhalten
\[
e(GX+HY)=\pm((\alpha g+\gamma h)X+(\beta g+\delta h)Y),\quad\text{\textit{d.h.}}
\]
\[
\pm eG=\alpha g+\gamma h,\quad\pm eH=\beta g+\delta h
\]
Bestimmen wir also $\mathfrak{G},\mathfrak{H}$ so, dass $\mathfrak{G}G+\mathfrak{H}H=1$ ist, so gilt
\[
\pm e=\mathfrak{G}(\alpha g+\gamma h)+\mathfrak{H}(\beta g+\delta h),\quad\text{also \textit{ganzzahlig.}\quad\textit{Q.E.D.}}
\]
Nehmen wir umgekehrt an, $\frac{M}{m}$ sei eine ganze Quadratzahl $=ee$, so enthält die Form $f$ die Form $F$. Wir können nämlich ganze Zahlen $\alpha,\beta,\gamma,\delta$ so bestimmen, dass
\SourceLeaf{218}
\[
\alpha g+\gamma h=\pm eG,\quad \beta g+\delta h=\pm eH
\]
gilt. Dazu wählen wir ganze Zahlen $\mathfrak{g},\mathfrak{h}$ so, dass $\mathfrak{g}g+\mathfrak{h}h=1$ ist. Jene Gleichungen werden erfüllt, wenn wir
\[
\begin{aligned}
\alpha&=\pm eG\mathfrak{g}+hz, &\quad\gamma&=\pm eG\mathfrak{h}-gz\\
\beta&=\pm eH\mathfrak{g}+hz', &\quad\delta&=\pm eH\mathfrak{h}-gz'
\end{aligned}
\]
setzen, welche ganzzahligen Werte wir auch $z,z'$ geben. Daher ist $F$ unter $f$ enthalten. \textit{Damit ist die Behauptung bewiesen.} Zugleich sehen wir leicht, dass diese Formeln alle möglichen Werte von $\alpha,\beta,\gamma,\delta$, \textit{d.h.} alle Transformationen von $f$ in $F$, angeben, wenn $z,z'$ alle ganzen Zahlen durchlaufen.

II. Für zwei gegebene Formen $f=axx+2bxy+cyy$ mit einer Determinante nicht $=0$ und $F=M(GX+HY)^2$ mit der Determinante $=0$ (wobei $G,H$ wie zuvor teilerfremd sind) behaupte ich \textit{erstens}: Enthält $f$ die Form $F$, so lässt sich die Zahl $M$ durch $f$ darstellen; \textit{zweitens}: Lässt sich $M$ durch $f$ darstellen, so ist $F$ unter $f$ enthalten; \textit{drittens}: Werden in diesem Fall alle Darstellungen von $M$ durch $f$ mit $x=\xi,y=\upsilon$ angegeben, so werden alle Transformationen von $f$ in $F$ durch $G\xi,H\zeta,G\upsilon,H\upsilon$ dargestellt. Dies alles beweisen wir folgendermaßen.

$1^{\mathrm{o}}$ Nehmen wir an, $f$ gehe in $F$ über, und zwar durch die Substitution $\alpha,\beta,\gamma,\delta$, und wählen wir $\mathfrak{G},\mathfrak{H}$ so, dass $\mathfrak{G}G+\mathfrak{H}H=1$ ist. Setzen wir dann $x=\alpha\mathfrak{G}+\beta\mathfrak{H}$, $y=\gamma\mathfrak{G}+\delta\mathfrak{H}$, so wird der Wert von $f$ offenbar gleich $M$, und somit ist $M$ durch $f$ darstellbar.

$2^{\mathrm{o}}$ Nehmen wir an, $a\xi\xi+2b\xi\upsilon+c\upsilon\upsilon=M$ sei erfüllt. Durch die Substitution $G\xi,H\xi,G\upsilon,H\upsilon$ geht die Form $f$ offenbar in $F$ über. Dass aber

$3^{\mathrm{o}}$ in diesem Fall die Substitution $G\xi,H\zeta,G\upsilon,H\upsilon$ \textit{alle} Transformationen von $f$ in $F$ angibt, wenn $\xi,\upsilon$ alle Werte von $x,y$ darstellen, für die $f=M$ ist, sehen wir so. Sei $\alpha,\beta,\gamma,\delta$ eine beliebige Transformation von $f$ in $F$, und wie zuvor gelte $\mathfrak{G}G+\mathfrak{H}H=1$. Unter den Werten von $x,y$ befinden sich dann auch
\[
x=\alpha\mathfrak{G}+\beta\mathfrak{H},\quad y=\gamma\mathfrak{G}+\delta\mathfrak{H}
\]
woraus sich die Substitution
\SourceLeaf{219}
\[
G(\alpha\mathfrak{G}+\beta\mathfrak{H}),\quad H(\alpha\mathfrak{G}+\beta\mathfrak{H}),\quad G(\gamma\mathfrak{G}+\delta\mathfrak{H}),\quad H(\gamma\mathfrak{G}+\delta\mathfrak{H})
\]
ergibt, oder
\[
\begin{aligned}
&\alpha+\mathfrak{H}(\beta G-\alpha H), &\quad&\beta+\mathfrak{G}(\alpha H-\beta G)\\
&\gamma+\mathfrak{H}(\delta G-\gamma H), &&\delta+\mathfrak{G}(\gamma H-\delta G)
\end{aligned}
\]
Da aber
\[
\begin{aligned}a(\alpha X+\beta Y)^2+2b(\alpha X+\beta Y)(\gamma X+\delta Y)&{}\\
{}+c(\gamma X+\delta Y)^2&=M(GX+HY)^2\end{aligned}
\]
gilt, ist
\[
\begin{aligned}
a(\alpha\delta-\beta\gamma)^2&=M(\delta G-\gamma H)^2\\
c(\beta\gamma-\alpha\delta)^2&=M(\beta G-\alpha H)^2
\end{aligned}
\]
und daher (weil die Determinante von $f$, mit $(\alpha\delta-\beta\gamma)^2$ multipliziert, gleich der Determinante von $F$, \textit{d.h.} $=0$, und somit auch $\alpha\delta-\beta\gamma=0$ ist)
\[
\delta G-\gamma H=0,\quad \beta G-\alpha H=0
\]
Jene Substitution geht damit in $\alpha,\beta,\gamma,\delta$ über. Daher liefert die angegebene Formel \textit{alle} Transformationen von $f$ in $F$.

III. Es bleibt zu zeigen, wie sich alle Darstellungen einer gegebenen Zahl durch eine gegebene Form mit der Determinante 0 angeben lassen. Sei die Form $m(gx+hy)^2$. Offenbar muss jene Zahl durch $m$ teilbar und der Quotient eine Quadratzahl sein. Setzen wir also die gegebene Zahl $=mee$, so ist für dieselben Werte von $x,y$, für die $m(gx+hy)^2=mee$ gilt, offenbar $gx+hy$ entweder $=+e$ oder $=-e$. Wir erhalten daher alle Darstellungen, indem wir alle ganzzahligen Lösungen der linearen Gleichungen $gx+hy=e$, $gx+hy=-e$ bestimmen. Dass diese lösbar sind, ist bekannt, denn $g,h$ sind wie vorausgesetzt teilerfremd. Bestimmen wir nämlich $\mathfrak{g},\mathfrak{h}$ so, dass $\mathfrak{g}g+\mathfrak{h}h=1$ ist, so erfüllen wir die erste Gleichung durch $x=\mathfrak{g}e+hz$, $y=\mathfrak{h}e-gz$, die zweite dagegen durch $x=-\mathfrak{g}e+hz$, $y=-\mathfrak{h}e-gz$, wobei $z$ eine beliebige ganze Zahl bezeichnet. Zugleich geben diese Formeln \textit{alle} ganzzahligen Werte von $x,y$ an, wenn $z$ alle ganzen Zahlen durchläuft.

